\chapter{Uncertainty}
In this chapter, the visualization of random variables and statistical data is discussed.
Uncertain data is variant data.
Uncertainty has to be visualized appropriately.
The problem is, that many completely different data sources can have similar or identical plots of them.
Therefore, identifying correlation between variables can not always be done by just comparing boxplots or directly plotting the data.
The takeaway is, that the importance lies in visualizing mathematical confidence, not intuitive confidence.

Visualization can again be done using certain chart idioms.
\begin{description}
    \item[Boxplot] The boxplot visualizes a variant data source, by dividing the vertical part of the plot into different segments.
        Outliers are marked as dots an channeled using the appropriate quantitative value.
        Then the minimum is displayed as a horizontal line, which is connected by a vertical line to the 25th percentile.
        The 25th percentile starts a rectangle which has another horizontal line inside of it, which is the median.
        The rectangle ends with the 75th percentile, which is connected by a vertical line with the maximum.
        The 25th and 75th percentiles are calculated, by splitting the dataset in two according to the mean value of the whole set and calculating the lower and higher mean values.

    \item[Gradient and Violin Plot] Gradient and violin plots are plots which encode the cumulative density of the area containing all values between minimum and maximum by transparency or width.
        For small datasets, these plots may confuse the reader and hallucinate meaning.
\end{description}

For bivariate data or temporal data, similar charts can be used.
